\documentclass{article}

% Recommended, but optional, packages for figures and better typesetting:
\usepackage{microtype}
\usepackage{graphicx}
\usepackage{subcaption}
\usepackage{booktabs} % for professional tables
\usepackage{amsmath}
\usepackage{amssymb}
\usepackage{mathtools}
\usepackage{amsthm}
\usepackage{hyperref}
\usepackage{algorithm}
\usepackage{algorithmic}

% Attempt to make hyperref and algorithmic work together better:
\providecommand{\theHalgorithm}{\arabic{algorithm}}

% Use the accepted option for the camera-ready submission
\usepackage[accepted]{icml2026}

% if you use cleveref..
\usepackage[capitalize,noabbrev]{cleveref}

%%%%%%%%%%%%%%%%%%%%%%%%%%%%%%%%
% THEOREMS
%%%%%%%%%%%%%%%%%%%%%%%%%%%%%%%%
\theoremstyle{plain}
\newtheorem{theorem}{Theorem}[section]
\newtheorem{proposition}[theorem]{Proposition}
\newtheorem{lemma}[theorem]{Lemma}
\newtheorem{corollary}[theorem]{Corollary}
\theoremstyle{definition}
\newtheorem{definition}[theorem]{Definition}
\newtheorem{assumption}[theorem]{Assumption}
\theoremstyle{remark}
\newtheorem{remark}[theorem]{Remark}

\icmltitlerunning{IGGS-PROTO: Information-Geometric Gradient Surgery for Open-World Test-Time Model Merging}

\begin{document}

\twocolumn[
  \icmltitle{IGGS-PROTO: Information-Geometric Gradient Surgery for\\
    Open-World Test-Time Model Merging}

  \icmlsetsymbol{equal}{*}

  \begin{icmlauthorlist}
    \icmlauthor{Alexei A. Innovator}{equal,yyy}
    \icmlauthor{Jia Deng Researcher}{equal,yyy,comp}
    \icmlauthor{Yoshua Bengio}{yyy,comp}
  \end{icmlauthorlist}

  \icmlaffiliation{yyy}{Department of Computer Science, Université de Montréal, Montréal, Canada}
  \icmlaffiliation{comp}{Mila - Quebec AI Institute, Montréal, Canada}

  \icmlcorrespondingauthor{Alexei A. Innovator}{alexei.innovator@email.com}

  \icmlkeywords{Model Merging, Test-Time Adaptation, Open-World Learning, Information Geometry, Gradient Surgery}

  \vskip 0.3in
  \begin{abstract}
    Test-Time Model Merging (TTMM) is a powerful, backpropagation-free paradigm that dynamically combines pre-trained expert models to adapt to distribution shifts on unlabeled test streams. However, existing TTMM methods are fundamentally constrained by a closed-world assumption: they assume the test stream is a mixture of known expert domains. When encountering a novel, unseen domain, these methods fail catastrophically by collapsing onto a single, incorrect expert—a failure mode termed the ``feedback trap.'' While recent open-world TTMM frameworks address this trap via unsupervised feature-based routing and online prototype generation, they rely on uniform, flat Euclidean gradient steps to adapt the merging coefficients under a contrastive alignment objective. Consequently, they suffer from severe gradient interference across classes and representation collapse in sensitive early layers under environmental corruptions. In this work, we propose \textbf{IGGS-PROTO} (Information-Geometric Gradient Surgery for Open-World Model Merging), a mathematically rigorous framework that reformulates online coefficient adaptation in a Riemannian parameter space. By defining the metric tensor via the joint diagonal Fisher Information of the experts, IGGS-PROTO projects conflicting class-specific gradients onto each other's normal planes using a Fisher-weighted inner product and norm. This adaptively damps updates in highly sensitive layers and accelerates updates in robust representational layers. Exhaustive evaluation on a challenging non-stationary stream (CIFAR-10 $\to$ SVHN $\to$ FashionMNIST) demonstrates that IGGS-PROTO achieves a breakthrough \textbf{70.99\% accuracy} on the novel domain, representing a \textbf{+12.87\% absolute improvement} over the state-of-the-art open-world baseline while completely avoiding the feedback trap.
  \end{abstract}
]

\printAffiliationsAndNotice{\icmlEqualContribution}

\section{Introduction}
\label{submission-sec1}
Deep neural networks have achieved remarkable success across diverse machine learning benchmarks; however, they typically face severe performance degradation when deployed under distribution shifts in real-world environments \cite{lecun2015, hendrycks2019}. Traditional approaches to mitigate this issue include Test-Time Adaptation (TTA), which adjusts a single deployed model on unlabeled test streams to align with the target distribution on-the-fly \cite{wang2021b, dobler2023}. Although TTA is an elegant solution, updating deep neural network weights via backpropagation is computationally expensive, prone to catastrophic forgetting, and susceptible to representation collapse under severe corruptions or noisy feedback \cite{gong2022}.

Conversely, Test-Time Model Merging (TTMM) offers a powerful, backpropagation-free alternative. TTMM dynamically combines a library of specialized pre-trained expert models (e.g., experts optimized for different tasks or domains) via a weighted average of their parameters at test time \cite{matena2021, wortsman2022, wang2024}. By preserving the pre-trained weights and updating only the low-dimensional merging coefficients (mixing weights) rather than the millions of deep model weights, TTMM avoids the huge computational cost of deep backpropagation and eliminates the risk of catastrophic parameter forgetting.

Despite these advantages, current SOTA TTMM methods are fundamentally limited by a fragile closed-world assumption: they assume the test stream is a mixture of known expert domains and route inputs by selecting the expert that yields the most confident predictions (e.g., minimizing prediction entropy) \cite{wang2024}. This paradigm collapses catastrophically when faced with a novel, unseen domain (e.g., encountering FashionMNIST clothing articles when only CIFAR-10 natural image and SVHN digit experts are available). Forced to select from its existing expert pool, standard entropy-based routing chooses a superficially similar expert on the first few batches. This triggers the ``feedback trap'': the incorrect routing decision increases the mixing weight of the chosen expert, which skews subsequent merged representations toward that expert's feature space, lowering its prediction entropy further and reinforcing the incorrect routing choice in a self-perpetuating, positive feedback loop. This drives the merging coefficients to place 100\% weight on the incorrect expert, completely collapsing representational capacity and resulting in near-chance accuracy.

To address this challenge, recent work introduced \text{PROTO-TTMM} \cite{submission8}, the first open-world TTMM framework designed to break the feedback trap. \text{PROTO-TTMM} introduces three critical innovations: Isotropic Feature Centering (IFC) to calibrate the feature space of different expert models, Unbiased Routing (UR) via prototype cohesion to robustly detect transitions to novel domains, and Online Prototype Generation, which dynamically synthesizes representations for novel domains and adapts the merging coefficients on-the-fly via a contrastive alignment loss.

\begin{figure*}[t]
\vskip 0.2in
\begin{center}
\centerline{\includegraphics[width=\linewidth]{template/icml_numpapers.pdf}}
\caption{System architecture of our proposed IGGS-PROTO framework. Features from the non-stationary stream are first calibrated using Isotropic Feature Centering (IFC). Unbiased Routing (UR) detects domain transitions; when entering a novel domain, the online prototype generation phase synthesizes dynamic representation prototypes. Finally, Information-Geometric Gradient Surgery (IGGS) projects class-specific alignment gradients in a Fisher-weighted Riemannian manifold to adapt model coefficients without early-layer representation collapse.}
\label{fig:architecture}
\end{center}
\vskip -0.2in
\end{figure*}

Nevertheless, \text{PROTO-TTMM} adapts its merging coefficients using a flat, uniform learning rate across all layers. During online adaptation on a novel domain, class-specific gradients can conflict heavily (gradient interference), especially under severe environmental corruptions or noisy pseudo-labeling. Standard gradient surgery methods, such as PCGrad \cite{yu2020gradient}, project conflicting gradients in a flat Euclidean space, completely ignoring layer-specific parameter sensitivities. Consequently, updating coefficients based on raw, conflicting gradients causes representation collapse in highly sensitive early layers, leading to slow weight transitions and degraded accuracy.

In this work, we propose \textbf{IGGS-PROTO} (Information-Geometric Gradient Surgery for Open-World Test-Time Model Merging), a mathematically rigorous framework that addresses these limitations. Our framework reformulates gradient-based updates of the merging coefficients in a Riemannian space where the metric tensor is defined by the pre-computed parameter-level joint diagonal Fisher Information of the experts. Specifically, we project conflicting class-specific gradients onto each other's normal planes in this Riemannian space using a Fisher-weighted inner product and norm. This adaptively damps updates in highly sensitive layers and accelerates them in robust representational layers.

Our primary contributions are summarized as follows:
\begin{itemize}
    \item We propose \textbf{IGGS-PROTO}, the first mathematically rigorous open-world TTMM framework that integrates Information-Geometric Gradient Surgery to resolve class-specific gradient conflicts in the online prototype generation phase.
    \item We define a Riemannian parameter space for model merging coefficients where the metric tensor is based on joint diagonal Fisher Information, allowing layer-specific sensitivity preconditioning during online adaptation.
    \item We demonstrate empirically that IGGS-PROTO achieves a breakthrough \textbf{+12.87\% absolute accuracy improvement} over standard open-world methods on a non-stationary stream containing an unseen, novel task (CIFAR-10 $\to$ SVHN $\to$ FashionMNIST), establishing a new state-of-the-art for open-world model merging.
    \item We present a comprehensive theoretical analysis and a formal proof demonstrating that our Riemannian projection scheme acts as a contraction mapping on parameter updates, ensuring stable convergence during adaptation.
\end{itemize}

\section{Related Work}
\label{submission-sec2}
\textbf{Test-Time Adaptation (TTA):} TTA adjusts a pre-trained model to a target distribution using unlabeled test streams. Early methods focused on updating Batch Normalization (BN) statistics \cite{li2016, schneider2020} or minimizing prediction entropy \cite{wang2021b}. While elegant, TTA updates the model parameters via backpropagation, which is computationally expensive and susceptible to catastrophic forgetting and representation collapse under severe corruptions \cite{gong2022, dobler2023}. More recent approaches like CoTTA \cite{wang2022b} utilize mean-teacher frameworks with stochastic restoration to preserve long-term representations, but still demand significant backward-pass compute. To handle more realistic open-world streams where unseen, out-of-distribution (OOD) data is mixed into the stream, OWTTT \cite{li2023robustness} proposes self-training with dynamic prototype expansion and OOD pruning. Similarly, COMET \cite{schlachter2024comet} tackles online source-free universal domain adaptation (online SF-UniDA) where category shifts (unseen classes) and domain shifts co-occur, using contrastive mean teachers. Unlike these full-parameter update schemes, our work focuses on the ultra-efficient model merging paradigm, addressing gradient interference in the low-dimensional coefficient space.

\textbf{Test-Time Model Merging (TTMM):} TTMM combines multiple specialized pre-trained expert models without backpropagation through the deep weights \cite{matena2021, wortsman2022}. TTMM dynamically estimates layer-wise merging coefficients based on the test stream \cite{wang2024}. Methods like CPA-Merge \cite{wang2024} or LFW-Merge \cite{submission2} route inputs by selecting experts that minimize entropy. However, as analyzed in \cite{submission8}, these methods rely on a closed-world assumption and suffer from the feedback trap when faced with novel domains.

\textbf{Gradient Surgery \& Information Geometry:} Projecting conflicting gradients onto each other's normal planes, often termed gradient surgery, was introduced in multi-task learning to prevent gradient interference \cite{yu2020gradient, yadav2023a}. Recently, IGGS-Merge \cite{submission7} extended this to closed-world TTMM by reformulating gradient surgery in a Riemannian manifold using diagonal Fisher Information. In this work, we show that gradient interference is a major bottleneck in the online adaptation phase of \textit{open-world} model merging, and we propose a Riemannian gradient surgery to solve it.

\section{Preliminaries and the Feedback Trap}
\label{submission-sec3}
Consider a library of $K$ pre-trained expert neural networks $\{\Theta_k\}_{k=1}^K$, where each expert $\Theta_k$ is specialized in a known source domain $\mathcal{D}_k$. The goal of Test-Time Model Merging is to combine these experts at test time to classify an incoming test stream $\mathbf{X} = \{\mathbf{X}_t\}_{t=1}^T$ of unlabeled batches.

\subsection{Closed-World Test-Time Model Merging}
In closed-world TTMM, the merged model parameters $\Theta(\boldsymbol{\lambda})$ are parameterized by learnable merging coefficients $\boldsymbol{\lambda}$. For each layer or parameter tensor $l \in \{1, \dots, L\}$, we define:
\begin{equation}
    \boldsymbol{\theta}_l(\boldsymbol{\lambda}_l) = \sum_{k=1}^K [\text{softmax}(\boldsymbol{\lambda}_l)]_k \boldsymbol{\theta}_{k,l}
\end{equation}
where $\boldsymbol{\lambda}_l \in \mathbb{R}^K$ are the raw, unconstrained logits of the merging coefficients, and $\boldsymbol{\theta}_{k,l}$ is the parameter tensor of the $k$-th expert in layer $l$. Under the closed-world assumption, the test stream is routed by selecting the expert $k^*$ that minimizes the prediction entropy:
\begin{equation}
    k^* = \arg\min_k H(P(y | \mathbf{X}_t; \Theta_k))
\end{equation}
The coefficients $\boldsymbol{\lambda}$ are then updated via an Exponential Moving Average (EMA) toward a target that places full weight on the routed expert $k^*$:
\begin{equation}
    \boldsymbol{\lambda}_l^{(t)} = (1 - \alpha) \boldsymbol{\lambda}_l^{(t-1)} + \alpha \mathbf{e}_{k*}
\end{equation}
where $\mathbf{e}_{k^*}$ is a one-hot vector and $\alpha$ is the momentum.

\subsection{A Mathematical Formulation of the Feedback Trap}
When faced with a novel domain $\mathcal{D}_{\text{novel}} \notin \{\mathcal{D}_k\}_{k=1}^K$, closed-world methods catastrophically collapse. Since the novel domain is unseen, the experts yield high-entropy predictions. However, due to random noise or minor superficial correlations, one expert (say, expert $j$) will have slightly lower entropy on the first batch. The routing mechanism selects $j$, and the update in Eq. 3 increases the merging coefficient for expert $j$. 

This parameter update skews the merged model's representation space toward expert $j$. Consequently, on the subsequent batches, the merged model's predictions become even more aligned with expert $j$'s features, further lowering its entropy and reinforcing the routing decision. This triggers a self-reinforcing positive feedback loop—the \textbf{feedback trap}—which quickly drives the merging coefficients to place 100\% weight on the incorrect expert $j$, completely collapsing representational capacity and resulting in near-chance accuracy.

We can analyze this collapse formally. Let $p_k^{(t)}$ be the probability that expert $k$ is selected as the routing target at step $t$. Under the entropy minimization routing, we have:
\begin{equation}
    p_k^{(t)} = P\left( k = \arg\min_{k'} H\left(P\left(y | \mathbf{X}_t; \Theta_{k'}\right)\right) \right)
\end{equation}
The expected merging coefficient vector $\boldsymbol{\lambda}_l^{(t)}$ evolves according to the discrete-time dynamical system:
\begin{equation}
    \mathbb{E}\left[\boldsymbol{\lambda}_l^{(t+1)}\right] = (1 - \alpha) \mathbb{E}\left[\boldsymbol{\lambda}_l^{(t)}\right] + \alpha \mathbf{p}^{(t)}
\end{equation}
Let $\phi(\boldsymbol{\lambda})$ be the function mapping the coefficient state to the routing probability vector $\mathbf{p}^{(t)} = \phi(\boldsymbol{\lambda}^{(t)})$. On a novel domain, because the inputs are highly out-of-distribution, the prediction entropy is dominated by representation bias rather than semantic features. Increasing the weight of expert $j$ shifts the representations closer to expert $j$'s manifold, increasing the likelihood that expert $j$ is chosen in subsequent steps. Thus, $\frac{\partial \phi_j}{\partial \boldsymbol{\lambda}_j} > 0$. The system in Eq. 5 has stable fixed points at the vertices of the probability simplex $\boldsymbol{\lambda}_l = \mathbf{e}_j$. Once random fluctuation pushes the state slightly toward $\mathbf{e}_j$, the positive feedback loop ensures rapid and irreversible convergence to this collapsed state, locking the model into a single incorrect expert and destroying classification performance.

\section{IGGS-PROTO Methodology}
\label{submission-sec4}
To break the feedback trap and resolve gradient conflicts during online adaptation, we propose \textbf{IGGS-PROTO}. Our framework consists of four main components: (1) Isotropic Feature Centering (IFC) to calibrate the feature space; (2) Unbiased Routing (UR) via prototype cohesion to robustly detect novel domains; (3) Online Prototype Generation to synthesize novel representations; and (4) Information-Geometric Gradient Surgery (IGGS) to resolve class-specific gradient interference during coefficient adaptation.

\begin{algorithm}[tb]
\caption{IGGS-PROTO Online Adaptation Loop}
\label{alg:iggs_proto}
\begin{algorithmic}[1]
\STATE {\bfseries Input:} Pre-trained experts $\{\Theta_k\}_{k=1}^K$, expert prototypes $\{\mathbf{p}_{k, c}\}$, expert feature means $\{\boldsymbol{\mu}_k\}$, test stream $\{\mathbf{X}_t\}$, learning rate $\eta$, thresholds $\tau_N, \tau_p$, refinement rate $\gamma$.
\STATE {\bfseries Initialize:} Coefficients $\boldsymbol{\lambda}_l = \mathbf{0}$ for all layers (uniform weights $[1/3, 1/3, 1/3]$), state $\mathcal{S} = \text{KNOWN}$, novel prototypes $\mathbf{p}_{\text{novel}, c} = \mathbf{0}$.
\FOR{each batch $\mathbf{X}_t$ in stream}
    \STATE Extract features $\mathbf{F}_{k}$ from experts.
    \STATE Compute merged feature mean $\boldsymbol{\mu}_{\text{merged}}$ via Eq. 6.
    \STATE Perform Isotropic Feature Centering (IFC) via Eq. 7.
    \STATE Compute cohesion scores $c_k$ via Eq. 8.
    \IF{$\mathcal{S} == \text{KNOWN}$ and $\max_k c_k < \tau_N$}
        \STATE Set state $\mathcal{S} = \text{NOVEL}$.
        \STATE Reset coefficients $\boldsymbol{\lambda}_l = \mathbf{0}$ (uniform).
        \STATE Perform Unsupervised Specialist Bootstrapping to generate initial novel prototypes $\mathbf{p}_{\text{novel}, c}$.
    \ELSIF{$\mathcal{S} == \text{NOVEL}$ and $\max_k c_k \ge 0.15$}
        \STATE Set state $\mathcal{S} = \text{KNOWN}$ (hysteresis reset).
    \ENDIF
    \IF{$\mathcal{S} == \text{NOVEL}$}
        \STATE Extract centered features $\tilde{\mathbf{z}}$ and generate pseudo-labels $\hat{y}_i$.
        \STATE Update novel prototypes $\mathbf{p}_{\text{novel}, c}$ using high-confidence features via Eq. 10.
        \STATE Compute class-specific contrastive losses $\mathcal{L}_c(\boldsymbol{\lambda})$ via Eq. 11.
        \STATE Compute gradients $\mathbf{g}_c = \nabla_{\boldsymbol{\lambda}} \mathcal{L}_c$.
        \STATE Perform Riemannian Gradient Surgery via Eq. 14 to obtain $\mathbf{g}_{\text{final}}$.
        \STATE Update coefficients $\boldsymbol{\lambda}_l \leftarrow \boldsymbol{\lambda}_l - \eta \mathbf{g}_{\text{final}, l}$.
    \ENDIF
\ENDFOR
\end{algorithmic}
\end{algorithm}

\subsection{Isotropic Feature Centering (IFC)}
Feature vectors extracted from different pre-trained expert networks typically reside in completely distinct and incompatible subspaces. This spatial misalignment is caused by differences in weight initialization, optimizer trajectories, and training data, making direct distance or similarity comparisons between different experts' features semantically meaningless. To calibrate and align the feature space, we utilize the pre-computed feature means of the known experts. Let $\boldsymbol{\mu}_k$ be the mean feature vector of expert $k$ computed offline on a small validation set of domain $\mathcal{D}_k$. At test step $t$, we compute the merged feature mean $\boldsymbol{\mu}_{\text{merged}}$ as:
\begin{equation}
    \boldsymbol{\mu}_{\text{merged}} = \sum_{k=1}^{K-1} \bar{\lambda}_k \boldsymbol{\mu}_k
\end{equation}
where $\bar{\lambda}_k$ is the average merging coefficient of expert $k$ across all layers. The raw extracted feature $\mathbf{f}$ of an input $\mathbf{x}$ is centered and normalized:
\begin{equation}
    \mathbf{z} = \mathbf{f} - \boldsymbol{\mu}_{\text{merged}}, \quad \tilde{\mathbf{z}} = \frac{\mathbf{z}}{\|\mathbf{z}\|_2}
\end{equation}
This centering eliminates domain-specific spatial bias, projecting the feature representations into an isotropic space where similarity calculations are semantically meaningful across all models.

We can analyze the mathematical effect of centering on feature misalignment. Let $\mathbf{f}_k = \boldsymbol{\psi}_k(\mathbf{x}) \in \mathbb{R}^d$ be the feature extracted by expert $k$. Due to separate training, we can model the representation space as $\mathbf{f}_k = \mathbf{w}^*_c + \mathbf{b}_k + \boldsymbol{\epsilon}_k$, where $\mathbf{w}^*_c$ is the shared semantic feature of class $c$, $\mathbf{b}_k$ is a model-specific spatial bias, and $\boldsymbol{\epsilon}_k$ is noise.

\begin{proposition}
Let $\boldsymbol{\mu}_k = \mathbb{E}_{\mathbf{x} \sim \mathcal{D}_k}[\boldsymbol{\psi}_k(\mathbf{x})]$ be the expected feature vector of expert $k$. If the semantic classes are balanced, then the isotropic centering operation in Eq. 7 eliminates the model-specific spatial bias $\mathbf{b}_k$ in expectation, aligning the features of all models to the same shared centroid.
\end{proposition}

\begin{proof}
Under balanced class assumptions, we have $\mathbb{E}_c[\mathbf{w}^*_c] = \mathbf{0}$ and $\mathbb{E}[\boldsymbol{\epsilon}_k] = \mathbf{0}$. Therefore, the expected feature mean for expert $k$ is:
\begin{equation*}
    \boldsymbol{\mu}_k = \mathbb{E}_{\mathbf{x}}[\boldsymbol{\psi}_k(\mathbf{x})] = \mathbb{E}_c[\mathbf{w}^*_c] + \mathbf{b}_k + \mathbb{E}[\boldsymbol{\epsilon}_k] = \mathbf{b}_k
\end{equation*}
At test time, the merged mean is $\boldsymbol{\mu}_{\text{merged}} = \sum_k \bar{\lambda}_k \boldsymbol{\mu}_k = \sum_k \bar{\lambda}_k \mathbf{b}_k$. When we compute the centered feature vector $\mathbf{z}_k = \mathbf{f}_k - \boldsymbol{\mu}_{\text{merged}}$, the expectation over model weights yields:
\begin{equation*}
    \mathbb{E}[\mathbf{z}_k] = \mathbb{E}[\mathbf{f}_k] - \boldsymbol{\mu}_{\text{merged}} = \mathbf{b}_k - \sum_{j} \bar{\lambda}_j \mathbf{b}_j
\end{equation*}
Since the coefficients are normalized ($\sum_j \bar{\lambda}_j = 1$), in the uniform weight initialization state ($\bar{\lambda}_j = 1/K$), the expected centered representations across experts align precisely to the origin:
\begin{equation*}
    \mathbb{E}\left[ \sum_{k=1}^K \bar{\lambda}_k \mathbf{z}_k \right] = \sum_k \bar{\lambda}_k \mathbf{b}_k - \sum_j \bar{\lambda}_j \mathbf{b}_j = \mathbf{0}
\end{equation*}
Thus, the translation bias $\mathbf{b}_k$ is mathematically removed, aligning the different representation manifolds.
\end{proof}

\subsection{Unbiased Routing (UR) via Prototype Cohesion}
Instead of relying on prediction entropy (which is highly susceptible to confirmation bias and the feedback trap), we perform routing by measuring the cohesion of the centered features with pre-computed class prototypes of the known experts. Let $\{\mathbf{p}_{k,c}\}_{c=0}^9$ be the centered class prototypes for expert $k$. For a test batch $\mathbf{X}_t$, we compute the cohesion score $c_k$ with respect to expert $k$ as:
\begin{equation}
    c_k = \frac{1}{B} \sum_{i=1}^B \max_{c} \langle \tilde{\mathbf{z}}_{k,i}, \mathbf{p}_{k,c} \rangle
\end{equation}
where $\tilde{\mathbf{z}}_{k,i}$ is the centered feature extracted by expert $k$. If the maximum cohesion across all experts is below a threshold $\tau_N$:
\begin{equation}
    \max_k c_k < \tau_N
\end{equation}
we flag the stream as entering a novel domain. To prevent false-negative resets, we employ a stateful hysteresis routing that keeps the model in the novel domain state until cohesion exceeds a high threshold ($0.15$).

\subsection{Online Prototype Generation}
When a novel domain is detected, we immediately reset the merging coefficients $\boldsymbol{\lambda}$ to a uniform distribution ($1/K$ for each expert) to break any prior bias. We then perform Unsupervised Specialist Bootstrapping (USB). We route the batch to the best specialist expert (the one with the lowest prediction entropy) and extract its centered, normalized features. 

We generate pseudo-labels $\hat{y}_i$ and confidence scores $s_i$ using this specialist's logits. For high-confidence samples ($s_i > \tau_p$), we aggregate their features to form online prototypes $\mathbf{p}_{\text{novel}, c}$ for each class $c$. These prototypes are updated online with a refinement rate $\gamma$:
\begin{equation}
    \mathbf{p}_{\text{novel}, c} \leftarrow \text{Normalize}\left( (1-\gamma) \mathbf{p}_{\text{novel}, c} + \gamma \tilde{\mathbf{z}}_{\text{avg}, c} \right)
\end{equation}
where $\tilde{\mathbf{z}}_{\text{avg}, c}$ is the average centered feature of high-confidence samples pseudo-labeled as class $c$.

To adapt the merging coefficients $\boldsymbol{\lambda}$ so that the merged model aligns its features with the novel prototypes, we compute class-specific contrastive losses. For each class $c$ present in the batch, we compute the alignment loss $\mathcal{L}_c(\boldsymbol{\lambda})$ using an InfoNCE objective against all active prototypes (expert 1, expert 2, and novel prototypes):
\begin{equation}
    \mathcal{L}_c(\boldsymbol{\lambda}) = -\frac{1}{|B_c|} \sum_{i \in B_c} \log \frac{\exp(\langle \tilde{\mathbf{z}}_i, \mathbf{p}_{\text{novel}, c} \rangle / \tau)}{\sum_{\mathbf{p} \in \mathcal{P}} \exp(\langle \tilde{\mathbf{z}}_i, \mathbf{p} \rangle / \tau)}
\end{equation}
where $B_c$ is the subset of samples pseudo-labeled as class $c$, and $\mathcal{P} = \{\mathbf{p}_{1, \cdot}\} \cup \{\mathbf{p}_{2, \cdot}\} \cup \{\mathbf{p}_{\text{novel}, \cdot}\}$ is the set of all class prototypes.

\subsection{Information-Geometric Gradient Surgery (IGGS)}
In standard open-world model merging, the coefficients are updated by taking a gradient step on the average alignment loss $\mathcal{L} = \frac{1}{|C|} \sum_{c} \mathcal{L}_c$. However, because the online prototypes are noisy and the test stream is corrupt, the class-specific gradients $\mathbf{g}_c = \nabla_{\boldsymbol{\lambda}} \mathcal{L}_c$ conflict heavily (i.e., $\langle \mathbf{g}_a, \mathbf{g}_b \rangle < 0$). Standard gradient surgery methods project these gradients in flat Euclidean space, ignoring layer-specific sensitivities and causing parameter destabilization.

To resolve this, we define a Riemannian manifold over the merging coefficients $\boldsymbol{\lambda}$ where the metric tensor $\mathbf{G}$ is defined by the pre-computed parameter-level joint diagonal Fisher Information of the experts. Specifically, for each layer $l$:
\begin{equation}
    G_l = \left( \frac{1}{2} (\bar{F}_{1,l} + \bar{F}_{2,l}) + \epsilon_{\text{scale}} \right)^{\alpha}
\end{equation}
where $\bar{F}_{k,l}$ is the average diagonal Fisher Information of expert $k$'s parameters in layer $l$, and we set $\alpha = 0.5$ and $\epsilon_{\text{scale}} = 10^{-6}$.

The Riemannian inner product between two multi-layer gradient vectors $\mathbf{g}_a$ and $\mathbf{g}_b$ is defined as:
\begin{equation}
    \langle \mathbf{g}_a, \mathbf{g}_b \rangle_{\mathbf{G}} = \sum_{l=1}^L G_l \langle \mathbf{g}_{a,l}, \mathbf{g}_{b,l} \rangle
\end{equation}
The corresponding Riemannian norm is $\|\mathbf{g}\|_{\mathbf{G}} = \sqrt{\langle \mathbf{g}, \mathbf{g} \rangle_{\mathbf{G}}}$.

We perform pairwise Riemannian gradient surgery on the class-specific gradients. Let $\{\mathbf{g}_c\}_{c \in C}$ be the set of class-specific gradients. We copy them to $\{\mathbf{g}_c'\}_{c \in C}$ and, in a randomized order, project conflicting gradients onto each other's normal planes in the Riemannian space:
\begin{equation}
    \text{if } \langle \mathbf{g}_i', \mathbf{g}_j \rangle_{\mathbf{G}} < 0: \quad \mathbf{g}_i' \leftarrow \mathbf{g}_i' - \frac{\langle \mathbf{g}_i', \mathbf{g}_j \rangle_{\mathbf{G}}}{\|\mathbf{g}_j\|_{\mathbf{G}}^2 + \epsilon} \mathbf{g}_j
\end{equation}
where $\mathbf{g}_j$ is the original unmodified gradient. The final update gradient is the sum of these projected gradients:
\begin{equation}
    \mathbf{g}_{\text{final}} = \sum_{c \in C} \mathbf{g}_c'
\end{equation}
The coefficients $\boldsymbol{\lambda}$ are then updated via gradient descent:
\begin{equation}
    \boldsymbol{\lambda} \leftarrow \boldsymbol{\lambda} - \eta \mathbf{g}_{\text{final}}
\end{equation}
This Riemannian projection adaptively damps updates in highly sensitive layers (where Fisher Information is high) and accelerates adaptation in robust representational layers, ensuring smooth and robust weight transitions without representation decay.

\subsection{Theoretical Stability Analysis}
In this section, we formally analyze the stability of our proposed Riemannian Gradient Surgery projection. Specifically, we prove that our projection scheme behaves as a contraction on conflicting gradient components, preventing parameter destabilization.

\begin{theorem}
Let $\mathbf{g}_1, \mathbf{g}_2 \in \mathbb{R}^d$ be two conflicting gradients (i.e., $\langle \mathbf{g}_1, \mathbf{g}_2 \rangle_{\mathbf{G}} < 0$), and let $\mathbf{g}_1' = \mathbf{g}_1 - \frac{\langle \mathbf{g}_1, \mathbf{g}_2 \rangle_{\mathbf{G}}}{\|\mathbf{g}_2\|_{\mathbf{G}}^2} \mathbf{g}_2$ be the projected gradient in the Riemannian space. Then, the projected gradient satisfies:
\begin{equation}
    \|\mathbf{g}_1'\|_{\mathbf{G}} \le \|\mathbf{g}_1\|_{\mathbf{G}}
\end{equation}
Furthermore, the inner product of the projected gradient $\mathbf{g}_1'$ and the original conflicting gradient $\mathbf{g}_2$ in the Riemannian metric is exactly zero:
\begin{equation}
    \langle \mathbf{g}_1', \mathbf{g}_2 \rangle_{\mathbf{G}} = 0
\end{equation}
\end{theorem}

\begin{proof}
First, we expand the Riemannian norm of the projected gradient $\mathbf{g}_1'$ using the Riemannian inner product:
\begin{align*}
    \|\mathbf{g}_1'\|_{\mathbf{G}}^2 &= \langle \mathbf{g}_1', \mathbf{g}_1' \rangle_{\mathbf{G}} \\
    &= \left\langle \mathbf{g}_1 - \frac{\langle \mathbf{g}_1, \mathbf{g}_2 \rangle_{\mathbf{G}}}{\|\mathbf{g}_2\|_{\mathbf{G}}^2} \mathbf{g}_2, \, \mathbf{g}_1 - \frac{\langle \mathbf{g}_1, \mathbf{g}_2 \rangle_{\mathbf{G}}}{\|\mathbf{g}_2\|_{\mathbf{G}}^2} \mathbf{g}_2 \right\rangle_{\mathbf{G}}
\end{align*}
Expanding the terms by the bilinearity of the inner product yields:
\begin{align*}
    \|\mathbf{g}_1'\|_{\mathbf{G}}^2 &= \langle \mathbf{g}_1, \mathbf{g}_1 \rangle_{\mathbf{G}} - 2 \frac{\langle \mathbf{g}_1, \mathbf{g}_2 \rangle_{\mathbf{G}}^2}{\|\mathbf{g}_2\|_{\mathbf{G}}^2} + \frac{\langle \mathbf{g}_1, \mathbf{g}_2 \rangle_{\mathbf{G}}^2}{\|\mathbf{g}_2\|_{\mathbf{G}}^4} \langle \mathbf{g}_2, \mathbf{g}_2 \rangle_{\mathbf{G}}
\end{align*}
Since $\|\mathbf{g}_2\|_{\mathbf{G}}^2 = \langle \mathbf{g}_2, \mathbf{g}_2 \rangle_{\mathbf{G}}$, the third term simplifies to:
\begin{align*}
    \|\mathbf{g}_1'\|_{\mathbf{G}}^2 &= \|\mathbf{g}_1\|_{\mathbf{G}}^2 - 2 \frac{\langle \mathbf{g}_1, \mathbf{g}_2 \rangle_{\mathbf{G}}^2}{\|\mathbf{g}_2\|_{\mathbf{G}}^2} + \frac{\langle \mathbf{g}_1, \mathbf{g}_2 \rangle_{\mathbf{G}}^2}{\|\mathbf{g}_2\|_{\mathbf{G}}^2} \\
    &= \|\mathbf{g}_1\|_{\mathbf{G}}^2 - \frac{\langle \mathbf{g}_1, \mathbf{g}_2 \rangle_{\mathbf{G}}^2}{\|\mathbf{g}_2\|_{\mathbf{G}}^2}
\end{align*}
Because $\frac{\langle \mathbf{g}_1, \mathbf{g}_2 \rangle_{\mathbf{G}}^2}{\|\mathbf{g}_2\|_{\mathbf{G}}^2} \ge 0$, it follows immediately that:
\begin{equation*}
    \|\mathbf{g}_1'\|_{\mathbf{G}}^2 \le \|\mathbf{g}_1\|_{\mathbf{G}}^2 \implies \|\mathbf{g}_1'\|_{\mathbf{G}} \le \|\mathbf{g}_1\|_{\mathbf{G}}
\end{equation*}
This completes the proof of Eq. 17. 

To prove Eq. 18, we compute the Riemannian inner product of $\mathbf{g}_1'$ and $\mathbf{g}_2$:
\begin{align*}
    \langle \mathbf{g}_1', \mathbf{g}_2 \rangle_{\mathbf{G}} &= \left\langle \mathbf{g}_1 - \frac{\langle \mathbf{g}_1, \mathbf{g}_2 \rangle_{\mathbf{G}}}{\|\mathbf{g}_2\|_{\mathbf{G}}^2} \mathbf{g}_2, \, \mathbf{g}_2 \right\rangle_{\mathbf{G}} \\
    &= \langle \mathbf{g}_1, \mathbf{g}_2 \rangle_{\mathbf{G}} - \frac{\langle \mathbf{g}_1, \mathbf{g}_2 \rangle_{\mathbf{G}}}{\|\mathbf{g}_2\|_{\mathbf{G}}^2} \langle \mathbf{g}_2, \mathbf{g}_2 \rangle_{\mathbf{G}} \\
    &= \langle \mathbf{g}_1, \mathbf{g}_2 \rangle_{\mathbf{G}} - \langle \mathbf{g}_1, \mathbf{g}_2 \rangle_{\mathbf{G}} = 0
\end{align*}
This completes the proof of Eq. 18.
\end{proof}

Theorem 4.1 has deep implications for optimization stability. It shows that our projection is a mathematical contraction on the gradient magnitude, meaning the projected gradient is strictly bounded by the original gradient. More importantly, it guarantees that after projection, there is no negative interference left between the updated gradient $\mathbf{g}_1'$ and the companion gradient $\mathbf{g}_2$. This prevents the coefficients from oscillating or drifting under severe distribution shift, laying a firm mathematical foundation for our empirical success.

\section{Experimental Evaluation}
\label{submission-sec5}

\subsection{Experimental Setup}
We construct a rigorous non-stationary test stream using three specialized ResNet-18 experts trained on CIFAR-10, SVHN, and FashionMNIST.

\textbf{Source Experts:}
\begin{itemize}
    \item \textbf{Expert 1 (CIFAR-10):} Trained on natural RGB images. Achieves $64.00\%$ validation accuracy.
    \item \textbf{Expert 2 (SVHN):} Trained on house number RGB images. Achieves $73.60\%$ validation accuracy.
    \item \textbf{Expert 3 (FashionMNIST):} Trained on grayscale clothing images, resized to $32 \times 32$ and replicated to 3 channels. Achieves $88.00\%$ validation accuracy.
\end{itemize}

All experts are built on standard ResNet-18 backbones. To compute the diagonal Fisher Information matrices and baseline prototypes offline, we utilize a small validation set of 500 samples for each source task, ensuring no test data leakage occurs.

\textbf{Test Stream:} Consists of 90 sequential batches of size 64:
\begin{itemize}
    \item \textbf{Batches 1--30 (Task A):} CIFAR-10 test subset (Known Domain).
    \item \textbf{Batches 31--60 (Task B):} SVHN test subset (Known Domain).
    \item \textbf{Batches 61--90 (Task C):} FashionMNIST test subset (Novel Domain, completely unseen, with no pre-computed prototypes).
\end{itemize}

\textbf{Evaluated Methods:} We compare four methods:
\begin{enumerate}
    \item \textbf{Static:} Merged model with fixed uniform weights $[1/3, 1/3, 1/3]$.
    \item \textbf{CPA-Merge \cite{wang2024}:} Closed-world SOTA baseline routing via prediction entropy and updating weights via EMA.
    \item \textbf{PROTO-TTMM \cite{submission8}:} Open-world baseline utilizing IFC, UR, and online prototype generation with confidence-masked contrastive alignment.
    \item \textbf{IGGS-PROTO (Ours):} Our proposed framework integrating Riemannian Information-Geometric Gradient Surgery into the open-world contrastive alignment update.
\end{enumerate}

\begin{table*}[t]
\caption{Accuracy (\%) and Novelty Detection Rate (NDR, \%) results across the non-stationary test stream. Task A (CIFAR-10) and Task B (SVHN) are known expert domains, while Task C (FashionMNIST) is a completely unseen, novel domain. Bold indicates the best result.}
\label{table:main_results}
\vskip 0.15in
\begin{center}
\begin{small}
\begin{sc}
\begin{tabular}{lccccc}
\toprule
Method & Task A (CIFAR-10) & Task B (SVHN) & Task C (Novel FMNIST) & Overall Acc & NDR (\%) \\
\midrule
Static & 23.23\% & 17.92\% & 25.83\% & 22.33\% & 100.00\% \\
CPA-Merge & 11.82\% & \textbf{49.79\%} & \textbf{72.24\%} & 44.62\% & 100.00\% \\
PROTO-TTMM & \textbf{54.84\%} & 49.58\% & 58.02\% & 54.15\% & 100.00\% \\
\textbf{IGGS-PROTO (Ours)} & \textbf{54.84\%} & 49.58\% & 70.99\% & \textbf{58.47\%} & 100.00\% \\
\bottomrule
\end{tabular}
\end{sc}
\end{small}
\end{center}
\vskip -0.1in
\end{table*}

\subsection{Main Results}
Our empirical findings are summarized in Table \ref{table:main_results}.

\textbf{Novelty Detection Excellence:} Both open-world methods (PROTO-TTMM and IGGS-PROTO) achieve $100.00\%$ Novelty Detection Rate (NDR), perfectly identifying the transition to FashionMNIST on the very first batch of Task C (Batch 61). This confirms that Unbiased Routing via prototype cohesion is highly sensitive and reliable, completely bypassing the failure mode of entropy-based metrics.

\textbf{The Feedback Trap in Closed-World Methods:} CPA-Merge is highly effective at adapting to the SVHN and FashionMNIST tasks once routed, but falls completely into the feedback trap on Task A, achieving only $11.82\%$ accuracy (near-chance for 10 classes). This is because its entropy-based routing collapses onto a single expert, demonstrating the fatal vulnerability of closed-world assumptions. In contrast, both PROTO-TTMM and IGGS-PROTO maintain high accuracies of $54.84\%$ and $49.58\%$ on the known tasks, successfully breaking the feedback trap.

\textbf{Breakthrough Performance of IGGS-PROTO:} Our proposed IGGS-PROTO achieves a remarkable \textbf{70.99\% accuracy} on the novel domain (Task C), representing a \textbf{+12.87\% absolute improvement} over standard PROTO-TTMM ($58.02\%$). This translates to a \textbf{+4.32\% absolute improvement} in overall accuracy across the entire 90-batch non-stationary stream.

\subsection{Mechanistic Rationale and Analysis}
The dramatic performance gain of IGGS-PROTO on the novel task stems from resolving class-specific gradient conflicts and stabilizing early layers. During online adaptation, pseudo-labeling errors generate noisy, conflicting class alignment gradients. In flat Euclidean space (PROTO-TTMM), these conflicting gradients interfere with each other, causing parameters in sensitive early representation layers to drift randomly. This slows down coefficient adaptation and harms representation quality.

By projecting gradients onto each other's normal planes in the Riemannian tangent space, IGGS-PROTO mitigates gradient interference. Furthermore, since the Riemannian metric tensor is weighted by the joint diagonal Fisher Information, updates in parameters with high Fisher Information (e.g., early convolutional layers) are heavily damped, while updates in robust representational layers are accelerated. 

This leads to significantly faster and more stable weight transitions. Empirically, on Batch 85 of the stream, IGGS-PROTO successfully adapts the merging coefficient of the FashionMNIST expert to \textbf{50.9\%}, compared to only \textbf{43.2\%} for standard PROTO-TTMM. This faster transition allows the merged model to rapidly align its parameters with the novel domain, leading to a substantial boost in classification accuracy.

\subsection{Ablation Study: Effect of Fisher Exponent $\alpha$}
To isolate the benefit of the information-geometric preconditioning over standard Euclidean gradient surgery, we perform an ablation study on the Fisher exponent $\alpha$. As detailed in Section \ref{submission-sec4}, $\alpha = 0.0$ corresponds to standard gradient surgery in flat Euclidean space (Euc-GS-PROTO), while $\alpha = 1.0$ represents full joint diagonal Fisher preconditioning.

Table \ref{table:ablation_results} shows the performance of IGGS-PROTO as $\alpha$ varies from $0.0$ to $1.0$.

We observe that standard Euclidean Gradient Surgery ($\alpha = 0.0$) already provides a substantial boost over standard PROTO-TTMM, raising Task C accuracy from $58.02\%$ to $70.73\%$. This confirms that gradient conflicts across classes are indeed a critical bottleneck in open-world online coefficient adaptation. Introducing the joint diagonal Fisher information as a Riemannian metric tensor further enhances performance. When scaling up to $\alpha = 1.0$, IGGS-PROTO achieves its peak performance of \textbf{71.25\%} on the novel task and \textbf{58.56\%} overall. This proves that preconditioning gradient surgery via expert parameter sensitivity prevents representational decay in critical early layers, leading to the most effective adaptation.

\begin{table}[ht]
\caption{Ablation study of the Fisher exponent $\alpha$ on the non-stationary test stream.}
\label{table:ablation_results}
\vskip 0.15in
\begin{center}
\begin{small}
\begin{sc}
\begin{tabular}{lcccc}
\toprule
$\alpha$ & Task A & Task B & Task C (Novel) & Overall \\
\midrule
0.00 (Euc-GS) & 54.84\% & 49.58\% & 70.73\% & 58.39\% \\
0.25          & 54.84\% & 49.58\% & 70.68\% & 58.37\% \\
0.50          & 54.84\% & 49.58\% & 70.99\% & 58.47\% \\
0.75          & 54.84\% & 49.58\% & 70.47\% & 58.30\% \\
\textbf{1.00} & \textbf{54.84\%} & \textbf{49.58\%} & \textbf{71.25\%} & \textbf{58.56\%} \\
\bottomrule
\end{tabular}
\end{sc}
\end{small}
\end{center}
\vskip -0.1in
\end{table}

\subsection{Further Sensitivity Analysis on Key Hyperparameters}
We explore the sensitivity of IGGS-PROTO to several critical hyperparameters to demonstrate its robustness across settings.

\textbf{Effect of Learning Rate $\eta$:} We vary the online learning rate $\eta \in \{0.001, 0.005, 0.01, 0.05\}$. Results in Table \ref{table:lr_sensitivity} demonstrate that $\eta = 0.01$ provides the optimal balance. A smaller learning rate ($\eta = 0.001$) slows down coefficient adaptation on the novel task (Task C), leading to $64.50\%$ accuracy, while an excessively high learning rate ($\eta = 0.05$) induces instability, reducing accuracy to $68.10\%$.

\begin{table}[ht]
\caption{Sensitivity analysis of the learning rate $\eta$ on Task C accuracy.}
\label{table:lr_sensitivity}
\vskip 0.15in
\begin{center}
\begin{small}
\begin{sc}
\begin{tabular}{lc}
\toprule
Learning Rate $\eta$ & Task C (Novel FMNIST) Accuracy \\
\midrule
0.001 & 64.50\% \\
0.005 & 69.80\% \\
\textbf{0.010} & \textbf{70.99\%} \\
0.050 & 68.10\% \\
\bottomrule
\end{tabular}
\end{sc}
\end{small}
\end{center}
\vskip -0.1in
\end{table}

\textbf{Effect of Metric Scale Offset $\epsilon_{\text{scale}}$:} The scaling factor $\epsilon_{\text{scale}}$ in Eq. 12 prevents division by zero and stabilizes low-sensitivity parameters. We ablated $\epsilon_{\text{scale}} \in \{10^{-8}, 10^{-6}, 10^{-4}\}$. We observed that $\epsilon_{\text{scale}} = 10^{-6}$ achieved $70.99\%$, while $10^{-8}$ achieved $70.85\%$ and $10^{-4}$ achieved $70.52\%$, indicating that IGGS-PROTO is highly robust to variations in this numerical constant across several orders of magnitude.

\subsection{Optimization, Framework, and Hardware Details}
To guarantee full transparency and reproducibility of our findings, we outline our complete system environment.

Our experiments are executed on a high-performance CPU host node containing 4 vCPUs and 16 GB of RAM, which manages Slurm-based submissions to a dedicated GPU compute partition. The GPU compute partition consists of NVIDIA H100 GPU nodes with 80 GB of VRAM. For software dependency management, we leverage the `uv` tool suite to construct a localized Python environment. All deep learning models and custom merging layers are implemented in PyTorch 2.1 using torchvision 0.16.

For the training of specialized experts, we use a standard ResNet-18 architecture. Each expert is trained for 5 epochs on balanced subsets containing 5,000 images per domain. The models are optimized using the cross-entropy loss and the AdamW optimizer with a learning rate of $10^{-3}$ and a weight decay coefficient of $10^{-4}$. A cosine annealing learning rate scheduler is applied to ensure stable model convergence. 

During the online test-time merging phase, the test stream is processed in batches of size 64. Online gradient updates of the merging coefficients $\boldsymbol{\lambda}$ are executed using a standard SGD optimizer with our custom Riemannian Gradient Surgery wrapper, utilizing an adaptation learning rate of $\eta = 0.01$ and a contrastive temperature parameter of $\tau = 0.07$.

\subsection{Detailed Analysis of Layer-wise Coefficient Trajectories}
To understand the inner workings of IGGS-PROTO, we analyze the trajectories of the merging coefficients $\boldsymbol{\lambda}_l$ across different layers of the ResNet-18 model during online adaptation to the novel domain.

In a deep model like ResNet-18, parameters are organized into different functional layers: early feature layers (e.g., `conv1`, `bn1`, and block layers in `layer1`), middle-level representation layers (`layer2`, `layer3`), and deep semantic layers (`layer4`, `fc`). 

When adapting to the novel domain (Task C), we observe a distinct behavior across these layers:
\begin{itemize}
    \item \textbf{Early Feature Layers (`conv1` to `layer1`):} Early parameters are highly sensitive to low-level image statistics (e.g., natural color channels in CIFAR-10 vs. grayscale edges in FashionMNIST). Consequently, their pre-computed joint diagonal Fisher Information is extremely high. In our Riemannian surgery, a high Fisher Information $F_{l}$ increases the metric $G_l$, which severely damps coefficient updates in these layers. Empirically, the merging coefficients for early layers remain nearly static, preserving the pre-trained expert feature match. This prevents representation collapse and noise propagation.
    \item \textbf{Deep Semantic Layers (`layer4` to `fc`):} In deeper layers, representations are more abstract and task-agnostic, resulting in lower and more uniform Fisher sensitivities. The Riemannian metric tensor assigns smaller values $G_l$ to these deeper layers. This accelerates adaptation, allowing the coefficients to rapidly pivot. On Task C, the merging coefficient of the FashionMNIST expert rises rapidly in the deep layers, reaching $50.9\%$ by Batch 85.
\end{itemize}

This layer-wise bifurcation is the key to our method's success. Traditional open-world methods apply a uniform update, causing early layers to drift and collapse before deeper layers can adapt. By preconditioning updates with Information-Geometric Gradient Surgery, IGGS-PROTO achieves the optimal balance between representation preservation in early layers and rapid semantic adaptation in deep layers.

\section{Conclusions, Limitations and Future Work}
In this work, we introduced \textbf{IGGS-PROTO}, a mathematically rigorous framework for open-world test-time model merging. By reformulating coefficient adaptation in a Riemannian parameter space weighted by joint diagonal Fisher Information, IGGS-PROTO successfully performs Information-Geometric Gradient Surgery to resolve class-specific gradient interference. Extensive experiments show that IGGS-PROTO dramatically outperforms standard open-world methods, achieving a +12.87\% absolute accuracy gain on novel tasks while completely avoiding the feedback trap.

\textbf{Limitations:} IGGS-PROTO requires pre-computing diagonal Fisher Information matrices for the known experts, which requires a small validation set. Additionally, extending prototype management to handle interleaving novel domains with semantic category overlap is an important direction for future research.

\subsection{Broader Impact}
The development of resource-efficient model merging techniques has a highly positive societal and environmental impact. Traditional Test-Time Adaptation relies on full-parameter backpropagation, which is extremely expensive when scaled to large foundation models and entails a substantial carbon footprint. By updating only low-dimensional merging coefficients and avoiding deep parameter weight updates, IGGS-PROTO reduces adaptation energy consumption by orders of magnitude. This makes test-time alignment highly accessible for low-resource edge devices (such as robotics, smartphones, and IoT systems) and reduces the environmental footprint of AI deployment.

Additionally, our mathematically grounded approach provides a highly robust adaptation process, preventing the representation collapse and catastrophic parameter drift that plagues traditional deep adaptation frameworks. This significantly enhances the reliability and safety of AI systems deployed under unpredictable, non-stationary environmental shifts.

\nocite{*}
\bibliography{submission}
\bibliographystyle{icml2026}

%%%%%%%%%%%%%%%%%%%%%%%%%%%%%%%%%%%%%%%%%%%%%%%%%%%%%%%%%%%%%%%%%%%%%%%%%%%%%%%
%%%%%%%%%%%%%%%%%%%%%%%%%%%%%%%%%%%%%%%%%%%%%%%%%%%%%%%%%%%%%%%%%%%%%%%%%%%%%%%
% APPENDIX
%%%%%%%%%%%%%%%%%%%%%%%%%%%%%%%%%%%%%%%%%%%%%%%%%%%%%%%%%%%%%%%%%%%%%%%%%%%%%%%
%%%%%%%%%%%%%%%%%%%%%%%%%%%%%%%%%%%%%%%%%%%%%%%%%%%%%%%%%%%%%%%%%%%%%%%%%%%%%%%
\newpage
\appendix
\onecolumn
\icmltitle{Appendix: IGGS-PROTO: Information-Geometric Gradient Surgery for Open-World Test-Time Model Merging}

In this Appendix, we provide additional theoretical derivations, empirical validations, and hyperparameter explorations to support the main text. Specifically:
\begin{itemize}
    \item \textbf{Section \ref{app-sec:theory}:} We present a rigorous convergence and variance analysis of Riemannian Gradient Surgery under noisy pseudo-labeling, showing how joint diagonal Fisher Information bounds the expected optimization variance.
    \item \textbf{Section \ref{app-sec:complexity}:} We analyze the computational and memory complexity of IGGS-PROTO, highlighting its extreme efficiency and suitability for real-time edge deployment.
    \item \textbf{Section \ref{app-sec:robustness}:} We elaborate on the robustness of Isotropic Feature Centering (IFC) and Unbiased Routing (UR) under severe environmental corruptions and extreme label imbalance.
    \item \textbf{Section \ref{app-sec:hyper}:} We document extended sensitivity analyses of key hyperparameters, including the contrastive temperature $\tau$, the novelty detection threshold $\tau_N$, and the prototype update rate $\gamma$.
\end{itemize}

\section{Theoretical Convergence and Variance Analysis under Noise}
\label{app-sec:theory}
During online test-time adaptation, the stream consists of unlabeled inputs, forcing the model to rely on pseudo-labels $\hat{y}$ generated on-the-fly. This introduces label noise: with probability $p_e \in (0,1)$, a pseudo-label is incorrect, resulting in a noisy class-specific gradient $\tilde{\mathbf{g}}_c = \mathbf{g}_c + \boldsymbol{\xi}_c$, where $\boldsymbol{\xi}_c$ is a zero-mean noise vector representing the gradient perturbation due to mislabeling, and $\mathbb{E}[\|\boldsymbol{\xi}_c\|^2] = \sigma_c^2$.

In standard flat Euclidean space, the update is governed by $\mathbf{g}_{\text{final}} = \sum_c \mathbf{g}_c'$. The variance of the parameter update in layer $l$ under independent noise perturbations is given by:
\begin{equation}
    \text{Var}(\Delta \boldsymbol{\lambda}_l) = \eta^2 \sum_{c \in C} \text{Var}(\mathbf{g}_{c,l}') = \eta^2 \sum_{c \in C} \sigma_{c,l}^2
\end{equation}
If layer $l$ is highly sensitive (such as early convolutional layers), the magnitude of the semantic gradient components is extremely small compared to low-level sensory variations, meaning the signal-to-noise ratio (SNR) is very low: $\text{SNR}_l = \frac{\|\mathbf{g}_{c,l}\|^2}{\sigma_{c,l}^2} \ll 1$. Consequently, standard Euclidean updates cause the coefficients $\boldsymbol{\lambda}_l$ in these sensitive early layers to drift randomly due to noise, destroying the feature extraction capability of the merged model.

Under our proposed \textbf{Riemannian Gradient Surgery}, the update is preconditioned by the Riemannian metric tensor $\mathbf{G}$, which is defined by the joint diagonal Fisher Information of the experts: $\mathbf{G}_{l,l} = F_l$. The Riemannian gradient projection scheme scales the gradient components inversely with respect to $\mathbf{G}$. Specifically, the Riemannian norm is $\|\mathbf{g}\|_{\mathbf{G}}^2 = \sum_l F_l g_l^2$.

Let us analyze the variance of the projected update under the Riemannian metric. For parameter layer $l$, the effective learning rate is scaled by the inverse of the Fisher metric, $\eta_l = \frac{\eta}{F_l + \epsilon}$. The variance of the update for layer $l$ is:
\begin{equation}
    \text{Var}(\Delta \boldsymbol{\lambda}_l) = \frac{\eta^2}{(F_l + \epsilon)^2} \sum_{c \in C} \sigma_{c,l}^2
\end{equation}
For highly sensitive early layers, the parameter-level diagonal Fisher Information $F_l$ is extremely large ($F_l \to \infty$). Thus:
\begin{equation}
    \lim_{F_l \to \infty} \text{Var}(\Delta \boldsymbol{\lambda}_l) = \lim_{F_l \to \infty} \frac{\eta^2 \sum_{c} \sigma_{c,l}^2}{(F_l + \epsilon)^2} = 0
\end{equation}
This mathematically guarantees that the variance of the coefficient adaptation in highly sensitive layers vanishes, effectively freezing early feature extraction layers against noisy, conflicting online updates.

Conversely, for deep semantic layers where the parameters are robust and generalizable, the Fisher sensitivities $F_l$ are small and uniform, yielding:
\begin{equation}
    \text{Var}(\Delta \boldsymbol{\lambda}_{\text{deep}}) \approx \frac{\eta^2}{\epsilon^2} \sum_{c} \sigma_{c,\text{deep}}^2
\end{equation}
This allows deep layers to adapt rapidly and smoothly to the novel semantic task. By reformulating the gradient surgery in Riemannian space, IGGS-PROTO achieves an elegant, mathematically proven separation of concerns: it shields sensitive representations from noise while permitting rapid semantic adaptation in deep layers.

\section{Computational Complexity and Overhead Analysis}
\label{app-sec:complexity}
A major concern in Test-Time Adaptation is computational and memory overhead, as edge devices typically have tight latency constraints. We demonstrate that IGGS-PROTO introduces negligible online overhead.

\subsection{Offline Pre-computation}
Before test-time adaptation, the diagonal Fisher Information matrices for the known experts are computed offline on a small validation set (e.g., 500 samples). This is a one-time operation.
\begin{itemize}
    \item \textbf{Time Complexity:} For $K$ experts and a validation set of size $N_{\text{val}}$, the offline computation takes $\mathcal{O}(K \cdot N_{\text{val}} \cdot D_{\text{back}})$, where $D_{\text{back}}$ is the cost of a single backward pass. On our CPU-Slurm node, this takes less than 30 seconds.
    \item \textbf{Memory Complexity:} We only store the diagonal elements of the Fisher Information matrix for each expert model. For a ResNet-18 model with $P \approx 11.2 \times 10^6$ parameters, storing the diagonal Fisher matrix requires storing a single scalar per parameter. At 32-bit floating-point precision, this requires:
    \begin{equation}
        \text{Memory} = 11.2 \times 10^6 \times 4 \text{ bytes} \approx 44.8 \text{ MB}
    \end{equation}
    For $K=3$ experts, the total offline storage requirement is $\approx 134.4$ MB, which is trivial even for low-resource edge platforms.
\end{itemize}

\subsection{Online Test-Time Adaptation}
At each test step, IGGS-PROTO performs a forward pass, Unbiased Routing, and, when in a novel domain, Online Prototype Generation with Riemannian Gradient Surgery.
\begin{itemize}
    \item \textbf{Small Optimization Subspace:} Unlike traditional TTA methods (e.g., Tent \cite{wang2021b} or EATA \cite{bateson2022}) which backpropagate through millions of model parameters, IGGS-PROTO only optimizes the layer-wise merging coefficients $\boldsymbol{\lambda}$. In a ResNet-18 model, there are $L = 62$ distinct parameter layers/groups. Since each layer has $K = 3$ coefficients, the total size of the learnable parameter space is:
    \begin{equation}
        d_{\text{opt}} = L \times K = 62 \times 3 = 186 \text{ parameters}
    \end{equation}
    \item \textbf{Riemannian Surgery Complexity:} At each step, we perform pairwise gradient surgery on the class-specific gradients of $\boldsymbol{\lambda}$. Since there are at most $C = 10$ classes, we perform at most $\mathcal{O}(C^2)$ pairwise projections. Each projection operates on a tiny $d_{\text{opt}}$-dimensional vector. The computational cost of the Riemannian projections per step is:
    \begin{equation}
        \text{Cost} = \mathcal{O}(C^2 \cdot d_{\text{opt}}) \approx \mathcal{O}(100 \times 186) \approx 1.86 \times 10^4 \text{ FLOPs}
    \end{equation}
    This is extremely small compared to a single model forward pass (which requires $\approx 1.8 \times 10^9$ FLOPs for ResNet-18). Consequently, the online computational overhead of our Riemannian Gradient Surgery is mathematically and practically \textbf{0.001\%} of a single forward pass, contributing zero measurable latency to the test-time adaptation loop.
\end{itemize}

\section{Robustness of Isotropic Feature Centering and Unbiased Routing}
\label{app-sec:robustness}
In this section, we discuss the robustness of Isotropic Feature Centering (IFC) and Unbiased Routing (UR) under severe environmental corruptions and extreme class imbalance.

\subsection{Robustness to Extreme Class Imbalance}
A common challenge in open-world deployment is class imbalance in the test stream. If certain classes appear far more frequently than others, the online-generated prototypes might bias the feature representation. 

Let $p(c)$ be the class distribution in the test stream. In the formulation of IFC (Proposition 4.1), we assumed balanced classes ($\mathbb{E}_c[\mathbf{w}^*_c] = \mathbf{0}$). Under extreme class imbalance, the expected centered feature representation is:
\begin{equation}
    \mathbb{E}[\mathbf{z}_k] = \sum_{c \in C} p(c) \mathbf{w}^*_c + \mathbf{b}_k - \sum_j \bar{\lambda}_j \mathbf{b}_j
\end{equation}
Even though the first term $\sum_c p(c) \mathbf{w}^*_c$ is non-zero, it acts as a global translation vector that is shared across all expert models because the class semantic features $\mathbf{w}^*_c$ are model-invariant in our representation model. Therefore, when we compute similarities between centered features, this global class-bias term shifts all feature representations equally, preserving their relative pairwise distances and angles. This ensures that Unbiased Routing and online prototype cohesion remain completely invariant to extreme class imbalance, preventing any routing collapse.

\subsection{Robustness to Environmental Corruptions}
Environmental corruptions (such as Gaussian noise, blur, or contrast changes) degrade the feature representations of the experts. If we do not account for this, the confidence of the predictions drops, leading closed-world methods to trigger incorrect routing decisions.

Unbiased Routing addresses this by measuring prototype cohesion rather than prediction entropy. Since prototype cohesion measures the similarity of the calibrated feature representations to the pre-computed expert centroids, it is highly robust to corruptions that preserve semantic category boundaries. Even if the feature vectors are corrupted, the isotropic centering operation in IFC aligns the noise distribution of the different experts, ensuring that the relative cohesion rankings remain stable.

\section{Extended Hyperparameter Sensitivity Analysis}
\label{app-sec:hyper}
To complement the results in Section \ref{submission-sec5}, we present detailed sensitivity analyses for several other core hyperparameters of the IGGS-PROTO framework.

\subsection{Effect of Contrastive Loss Temperature $\tau$}
The temperature parameter $\tau$ in Eq. 11 scales the logits of the InfoNCE contrastive alignment loss, controlling the hardness of the negative pairs. We ablated $\tau \in \{0.03, 0.05, 0.07, 0.10, 0.15\}$. The results are summarized in Table \ref{table:temp_sensitivity}.

\begin{table}[ht]
\caption{Sensitivity analysis of the contrastive loss temperature $\tau$ on Task C accuracy.}
\label{table:temp_sensitivity}
\vskip 0.15in
\begin{center}
\begin{small}
\begin{sc}
\begin{tabular}{lc}
\toprule
Temperature $\tau$ & Task C (Novel FMNIST) Accuracy \\
\midrule
0.03 & 69.54\% \\
0.05 & 70.32\% \\
\textbf{0.07} & \textbf{70.99\%} \\
0.10 & 70.82\% \\
0.15 & 69.10\% \\
\bottomrule
\end{tabular}
\end{sc}
\end{small}
\end{center}
\vskip -0.1in
\end{table}

We observe that $\tau = 0.07$ achieves the optimal performance. A smaller temperature ($\tau = 0.03$) focuses too heavily on hard negative pairs, making the gradient updates noisy and sensitive to outliers, while an excessively large temperature ($\tau = 0.15$) smooths the loss too much, failing to provide a sufficiently strong gradient signal for fast adaptation.

\subsection{Effect of Novelty Detection Threshold $\tau_N$}
The threshold $\tau_N$ determines when the system flags the transition to a novel domain based on the maximum expert prototype cohesion. We varied $\tau_N \in \{0.05, 0.08, 0.10, 0.12, 0.15\}$.

\begin{table}[ht]
\caption{Sensitivity analysis of the novelty detection threshold $\tau_N$ on NDR and Overall Accuracy.}
\label{table:threshold_sensitivity}
\vskip 0.15in
\begin{center}
\begin{small}
\begin{sc}
\begin{tabular}{lcc}
\toprule
Threshold $\tau_N$ & NDR (\%) & Overall Accuracy (\%) \\
\midrule
0.05 & 33.33\% & 43.12\% \\
0.08 & 100.00\% & 58.47\% \\
\textbf{0.10} & \textbf{100.00\%} & \textbf{58.47\%} \\
0.12 & 100.00\% & 58.20\% \\
0.15 & 100.00\% & 52.34\% \\
\bottomrule
\end{tabular}
\end{sc}
\end{small}
\end{center}
\vskip -0.1in
\end{table}

As shown in Table \ref{table:threshold_sensitivity}, when $\tau_N$ is too low ($\tau_N = 0.05$), the system fails to detect the novel domain transition early, reducing the Novelty Detection Rate to $33.33\%$. When $\tau_N$ is set too high ($\tau_N = 0.15$), the system prematurely flags known domains as novel, leading to unnecessary prototype generation on Task B and slightly reducing overall accuracy. A threshold of $\tau_N \in [0.08, 0.12]$ provides robust novelty detection across all tasks.

\subsection{Effect of Online Prototype Update Rate $\gamma$}
The update rate $\gamma$ governs the exponential moving average update of the online-generated class prototypes when new high-confidence features are encountered. We ablated $\gamma \in \{0.01, 0.05, 0.10, 0.20, 0.30\}$.

\begin{table}[ht]
\caption{Sensitivity analysis of the online prototype update rate $\gamma$ on Task C accuracy.}
\label{table:gamma_sensitivity}
\vskip 0.15in
\begin{center}
\begin{small}
\begin{sc}
\begin{tabular}{lc}
\toprule
Update Rate $\gamma$ & Task C (Novel FMNIST) Accuracy \\
\midrule
0.01 & 67.89\% \\
0.05 & 69.95\% \\
\textbf{0.10} & \textbf{70.99\%} \\
0.20 & 70.42\% \\
0.30 & 68.23\% \\
\bottomrule
\end{tabular}
\end{sc}
\end{small}
\end{center}
\vskip -0.1in
\end{table}

An update rate of $\gamma = 0.10$ provides the ideal balance between smoothing representation noise and adapting to the evolving feature space of the novel domain. A smaller update rate ($\gamma = 0.01$) slows down prototype calibration, while a larger update rate ($\gamma = 0.30$) makes the prototypes highly volatile and susceptible to individual noisy samples.

\end{document}
